\documentclass[10pt,a4paper]{amsart}
\usepackage[margin=19mm]{geometry}
\usepackage{amsmath,amssymb,mathtools}
\usepackage[scr=rsfs,cal=euler]{mathalfa}
\usepackage{xfrac}
\usepackage[unicode,hidelinks,bookmarks=false]{hyperref}
\newcommand{\ie}{i.e.\ }
\newcommand{\cf}{cf.\ }
\newcommand{\resp}{respectively\ }
\newcommand{\Subsection}{\subsection}
\providecommand{\nobreakdash}{\nobreak\hbox{-}}
\setlength{\emergencystretch}{2em}
\multlinegap0pt
\usepackage{bm}
\usepackage{bbm}
\usepackage{dsfont}
\usepackage{xspace}
\usepackage{cancel}
\usepackage{mathtools}
\usepackage{amsthm}
\makeatletter
\@ifundefined{theorem}{\newtheorem{theorem}{Theorem}}{}
\@ifundefined{corollary}{\newtheorem{corollary}[theorem]{Corollary}}{}
\@ifundefined{lemma}{\newtheorem{lemma}[theorem]{Lemma}}{}
\@ifundefined{proposition}{\newtheorem{proposition}[theorem]{Proposition}}{}
\makeatother
\title[Super Hayashi Quandles]{Super Hayashi Quandles}
\author{Antonio Lages, Pedro Lopes}
\date{Source: arXiv:2405.12329v2 (CC BY 4.0)}
\begin{document}
\maketitle

\noindent\emph{Source and licence.} Super Hayashi Quandles. Version arXiv:2405.12329v2 (CC BY 4.0), distributed under the Creative Commons Attribution 4.0 International licence, as recorded in the arXiv licence metadata for this exact version and shown on the abstract page https://arxiv.org/abs/2405.12329v2. Full rights notice: licenses/arxiv-2405.12329-CC-BY-4.0.txt.\par\smallskip

\noindent\textit{Editorial scope.} Counted unit: the Lemma whose source label is \texttt{cor:final} in the article (source0.tex lines 471--473, source label \texttt{cor:final}), delivered together with the article's own complete original proof of it (source0.tex lines 475--506, one \texttt{proof} environment of 32 lines). No mathematics is written, repaired or re-derived: statement and proof are the article's own text line for line. Required context retained uncounted: prop:assertions (lines 226--233); lem:Fxiso (lines 366--369); lem:kappa (lines 384--386); prop:asspowers (lines 419--421); cor:Bni+1 (lines 460--462). Every definition, display and label of those lines is retained unchanged. Omitted from this package and not redistributed: 1--225, 234--365, 370--383, 387--418, 422--459, 463--470, 474--474, 507--754. Nothing inside a retained excerpt is omitted or altered: every retained body is delivered exactly as the article's lines are written, so no omission receipt of any kind accompanies this package. Citations: the retained excerpts carry no citation command at all, so no bibliography entry of the article is transcribed and no reference is redistributed. Reference adaptation: none was needed and none was made; every reference inside the delivered unit resolves inside it, so no unresolved reference or citation is produced. Printed slips kept literal: no printed word, symbol, hypothesis or number of the counted statement or of its proof was altered. Layout and class: the article's own class is \texttt{article} with packages including amsmath, amssymb, fullpage, amsthm, xcolor, enumerate, ulem, caption, indentfirst; the wrapper is the lane's pinned \texttt{amsart} editorial wrapper at 10pt on A4, which restates the theorem-like environments the retained text needs and adds the article's own macro definitions that the retained text needs (none), with extra wrapper packages (bm, bbm, dsfont, xspace, cancel, mathtools). The delivered numbering is the unit's own. Assets: no retained span contains a figure environment, a graphics-inclusion command, a PDF-page inclusion, a table or a TikZ picture; no image file is copied into this package, the package declares no asset dependency and no image file is redistributed with it. Licence: version arXiv:2405.12329v2 of this article is distributed under the Creative Commons Attribution 4.0 International licence, as recorded in the arXiv licence metadata for this exact version and shown on the abstract page https://arxiv.org/abs/2405.12329v2. A full rights notice is delivered at licenses/arxiv-2405.12329-CC-BY-4.0.txt. Characters: every retained excerpt is pure ASCII, so the delivered engine, XeLaTeX with its default Latin Modern text font, reports no missing character.
\begin{proposition}\label{prop:assertions}
Let $(X, *)$ be a finite quandle.
Modulo isomorphism, the right multiplications of $(X,*)$ satisfy the following conditions:
\begin{enumerate}
    \item $R_1=(n_1)(n_2-\ell_2+1\quad n_2-\ell_2+2\quad\cdots\quad n_2)\cdots(n_c-\ell_c+1\quad n_c-\ell_c+2\quad\cdots\quad n_c)$;
    \item $R_{n_{i-1}+k}=R_1^{k}R_{n_i}R_1^{-k}$, for each $k\in\{1,\dots,\ell_i\}$ and $i\in\{2,\dots,c\}$.
\end{enumerate}
\end{proposition}

\begin{lemma}\label{lem:Fxiso}
For any $x\in X_{i+1}$,
$$(X_i, *) = (F_1, *) \cong (F_x, *)\qquad \text{ and } \qquad n_i=|F_x| .$$
\end{lemma}

\begin{lemma}\label{lem:kappa}
$$\ell_{i+1}=\kappa \ell(\ell+1)^{i-1} ,\qquad \text{ for some $\kappa\in\mathbf{Z}^+$}.$$
\end{lemma}

\begin{proposition}\label{prop:asspowers}
If $x,y\in X_{i+1}$ are such that $y\in B_x$, then $R_y^\ell=R_x^\ell$.
\end{proposition}

\begin{corollary}\label{cor:Bni+1}
$(B_{n_{i+1}},  \dot{+})$ is a $($cyclic$)$ subgroup of $(L_{i+1}, \dot{+})$, Then, the elements of $B_{n_{i+1}}$ are evenly spaced within $L_{i+1}$: $$B_{n_{i+1}}=\{n_i+j\kappa \ell (\ell+1)^{i-2}\, |\, j\in\{ 1, 2, \dots, \ell+1\}\} .$$
\end{corollary}

\begin{lemma}\label{cor:final}
$\ell_{i+1}=\ell\cdot(\ell+1)^{i-1}$.
\end{lemma}

\begin{proof}

Since $|X_{i+1}|=n_{i+1}=n_i+\ell_{i+1}$, $|F_{n_{i+1}}|=n_i$ (Lemma \ref{lem:Fxiso}) and $F_{n_{i+1}}\subset X_{i+1}$, then $|X_{i+1}\setminus F_{n_{i+1}}|=\ell_{i+1}$. Let $X_{i+1}\setminus F_{n_{i+1}}=\{g_1,g_2,\dots,g_{\ell_{i+1}}\}$. Without loss of generality,
\[R_{n_{i+1}}|_{X_{i+1}\setminus F_{n_{i+1}}}=(g_1\quad g_2\quad\dots\quad g_{\ell_{i+1}}) ,\]since $F_{n_{i+1}}$ is the union of the cycles of $R_{n_ {i+1}}$ whose lengths are less than or equal to $\ell_i$.
By assertion \textit{2.} in Proposition \ref{prop:assertions}, $R_{n_i+j\cdot\kappa\cdot\ell\cdot(\ell+1)^{i-2}}=R_1^{j\cdot\kappa\cdot\ell\cdot(\ell+1)^{i-2}}R_{n_{i+1}}R_1^{-j\cdot\kappa\cdot\ell\cdot(\ell+1)^{i-2}}$, $\forall j\in\{1,\dots,\ell+1\}$, so,
\[R_{n_i+j\cdot\kappa\cdot\ell\cdot(\ell+1)^{i-2}}|_{X_{i+1}\setminus F_{n_{i+1}}}=(R_1^{j\cdot\kappa\cdot\ell\cdot(\ell+1)^{i-2}}(g_1)\quad R_1^{j\cdot\kappa\cdot\ell\cdot(\ell+1)^{i-2}}(g_2)\quad\cdots\quad R_1^{j\cdot\kappa\cdot\ell\cdot(\ell+1)^{i-2}}(g_{\ell_{i+1}})),\]
and also, its $l$-th iterate,
\[R_{n_i+j\cdot\kappa\cdot\ell\cdot(\ell+1)^{i-2}}^\ell|_{X_{i+1}\setminus F_{n_{i+1}}}=(R_1^{j\cdot\kappa\cdot\ell\cdot(\ell+1)^{i-2}}(g_1)\quad R_1^{j\cdot\kappa\cdot\ell\cdot(\ell+1)^{i-2}}(g_{\ell+1})\quad\cdots\quad R_1^{j\cdot\kappa\cdot\ell\cdot(\ell+1)^{i-2}}(g_{\ell_{i+1}-\ell+1}))\]
\[(R_1^{j\cdot\kappa\cdot\ell\cdot(\ell+1)^{i-2}}(g_2)\quad R_1^{j\cdot\kappa\cdot\ell\cdot(\ell+1)^{i-2}}(g_{\ell+2})\quad\cdots\quad R_1^{j\cdot\kappa\cdot\ell\cdot(\ell+1)^{i-2}}(g_{\ell_{i+1}-\ell+2}))\cdots\]
\[\cdots(R_1^{j\cdot\kappa\cdot\ell\cdot(\ell+1)^{i-2}}(g_\ell)\quad R_1^{j\cdot\kappa\cdot\ell\cdot(\ell+1)^{i-2}}(g_{2\cdot\ell})\quad\cdots\quad R_1^{j\cdot\kappa\cdot\ell\cdot(\ell+1)^{i-2}}(g_{\ell_{i+1}})).\]
Since, for each $j\in \{ 1, 2, \dots , \ell +1 \}$, $$n_i+j\cdot\kappa\cdot\ell\cdot(\ell+1)^{i-2}\in B_{n_{i+1}} \qquad \qquad   (\text{Corollary \ref{cor:Bni+1}}) ,$$ then  $$R_{n_i+j\cdot\kappa\cdot\ell\cdot(\ell+1)^{i-2}}^\ell$$ does not depend on $j$ (Proposition \ref{prop:asspowers}). Since $1\in F_1$ and the $F_k$'s partition $X_{i+1}$, $1\in X_{i+1}\setminus F_{n_{i+1}}$ and we may assume that $g_\ell=1$. Thus,
\begin{align*}
(R_1^{j\cdot\kappa\cdot\ell\cdot(\ell+1)^{i-2}}(g_\ell)&\quad R_1^{j\cdot\kappa\cdot\ell\cdot (\ell+1)^{i-2}}(g_{2\cdot\ell})\quad\cdots\quad R_1^{j\cdot\kappa\cdot\ell\cdot(\ell+1)^{i-2}}(g_{\ell_{i+1}}))=\\
&=(1\quad R_1^{j\cdot\kappa\cdot\ell\cdot(\ell+1)^{i-2}}(g_{2\cdot\ell})\quad\cdots\quad R_1^{j\cdot\kappa\cdot\ell\cdot(\ell+1)^{i-2}}(g_{\ell_{i+1}}))=(1\quad g_{2\cdot\ell}'\quad g_{3\cdot\ell}'\quad\cdots\quad g_{\ell_{i+1}}')
\end{align*} does not depend on $j$.

\begin{lemma}\label{lem:GF1}
$$\{ g_{\ell}, g_{2\ell}, g_{3\ell}, \cdots  , g_{\ell_{i+1}}  \} = \{ 1, g_{2\ell}, g_{3\ell}, \cdots , g_{[\ell_{i+1}/\ell]\ell}\}\subset F_1 .$$
\end{lemma}
\begin{proof}
Assume to the contrary and suppose the inclusion is not true. Since $1\in F_1$, there exists $s\in \{ 2, \dots , \ell_{i+1}/\ell \}$ such that $$g_{s\ell}\notin F_1=\{ y\in X_{i+1}\,|\, R_1^{\ell_i}(y)=y \} .$$ Then, $g_{s\ell}\in X_{i+1}\setminus F_1$. Specifically, $g_{s\ell}$ belongs to a disjoint cycle of $R_1$ of length greater than $\ell_i$ (since the elements in $F_1$ belong to cycles whose lengths divide $\ell_i$). Since we are working within the restriction to $X_{i+1}$, there is only one such cycle. It is the cycle of length $l_{i+1}=\kappa \ell (\ell+1)^{i-1}$ (Lemma \ref{lem:kappa}). Since $j\kappa \ell(\ell +1)^{i-2}< \kappa\ell(\ell +1)^{i-1}$, for $1\leq j \leq \ell$, $$|\{ R_1^{j\kappa \ell(\ell +1)^{i-2}}(g_{s\ell})\, |\, 1\leq j \leq \ell \}|=\ell .$$  Therefore, the sequence of permutations $$\Bigg(\quad  \bigg(1\quad R_1^{j\cdot\kappa\cdot\ell\cdot(\ell+1)^{i-2}}(g_{2\cdot\ell})\quad\cdots\quad R_1^{j\kappa \ell(\ell +1)^{i-2}}(g_{s\ell})\quad\cdots\quad R_1^{j\cdot\kappa\cdot\ell\cdot(\ell+1)^{i-2}}(g_{\ell_{i+1}})\bigg)\quad |\quad 1\leq j \leq \ell+1 \quad\Bigg) $$ is not constant which contradicts Proposition \ref{prop:asspowers}. The proof of Lemma \ref{lem:GF1} is complete.

\end{proof}
We now resume the proof of Lemma \ref{cor:final}: \begin{align*}
&\{1,g_{2\cdot\ell},g_{3\cdot\ell},\dots,g_{(\ell_{i+1}/\ell)\ell}\}\subseteq F_1 \quad \Longrightarrow \quad |\{1,g_{2\cdot\ell},g_{3\cdot\ell},\dots,g_{(\ell_{i+1}/\ell)\ell}\}|\leq |F_1| \quad \Longrightarrow \quad \ell_{i+1}/\ell \leq n_i \quad \Longleftrightarrow \quad \\
&\kappa \ell(\ell+1)^{i-1}/\ell \leq (\ell+1)^{i-1} \quad \Longrightarrow \quad \kappa = 1 \qquad \text{ since }\kappa\in\mathbf{Z}^+
\end{align*}
where the equivalence above follows from Lemma \ref{lem:kappa} and from the induction hypothesis. The proof of Lemma \ref{cor:final} is complete.



\end{proof}

\end{document}
